\section{xv6调度机制分析}
本章对xv6操作系统中的进程调度机制进行分析。首先介绍xv6进程模型与调度机制，然后分析其默认的时间片轮转调度算法及其特点，最后从决策的角度对任务调度进行形式化建模，为后续基于强化学习的调度策略设计提供理论基础。
\subsection{xv6进程调度机制}
\subsubsection{进程模型}
进程是调度原子。
进程是状态机。
进程是操作系统中的信息结构体（或者说数据结构？）
\subsubsection{调度循环}
1. 循环结构 -> 长期的决策结构
2. 获取状态（扫描查找proc[]） -> 选择策略（选择进程） ->  目的是为了完成进程 + 增加CPU的使用率（是否有中断睡眠）-> 调度完之后，会改变proc[]的状态（比如有进程进入SLEEPING状态，ZOMBIE状态）
3. 长期序贯决策，收益最大化的结构，调度机制。
\subsubsection{上下文切换}
决策与线程执行的关键机制。
\subsection{xv6默认算法分析}
\subsubsection{轮转调度原理}
\subsubsection{优缺点分析}
基于固定规则的静态调度策略，
不依赖于系统状态，
不利用历史信息
策略不发生变化
\subsection{调度问题的结构建模}
根据xv6本身的调度机制特点
长期序贯决策，
利用系统状态信息，利用历史信息，能够使用更加合适的策略进行决策
所以使用MDP进行形式化建模